% Source: https://artofproblemsolving.com/community/c3h3242498_the_polymath_competition
% Stable identity: aops-topic-3242498-q22
An even integer with fewer than five divisors must be $2,4,8$, or $2q$ with $q$ an odd prime, by the divisor-count formula. Since $\phi(p)=p-1$ for prime $p$, the initial possibilities are $p=2,3,5$ or $p=2q+1$ with odd prime $q$; $p-1=8$ gives the composite $9$.
In the latter case $\phi(2q)=q-1=\phi(q)$, so $p$ works exactly when the smaller prime $q$ works. Starting with $2,3,5$ and repeatedly testing $2q+1$ therefore gives the complete list
\[2,\ 3,\ 5,\ 7,\ 11,\ 23,\ 47.\]
The branches stop because $2(7)+1=15$ and $2(47)+1=95$ are composite; the other branches link back into this list. Each listed prime indeed works, since its iterates follow this recursion and eventually reach $4,2,1$. Their sum is $98$, with digit sum $\boxed{17}$, choice E.
